\chapter{Formal Model, Research Problem, and Assumptions}
\label{chap:model}

\section{Network Configuration}
\label{sec:model:network}

\TODO{Define the physical parameters: number of hops $N$, per-photon
loss and gate error rates, memory coherence time, and the half-RGS
branching factor. Explain which parameters are fixed in this work and
which are swept in experiments.}

\section{State Space}
\label{sec:model:state}

Each edge is represented by a state object
\[
    S = \bigl(\mathbf{e},\; s,\; t,\; t_{\mathrm{gen}},\; \kappa,\; r,\; h\bigr),
\]
where $\mathbf{e}$ is the error vector, $s$ the accumulated Z side-effect
parity, $t$ and $t_{\mathrm{gen}}$ the current and generation times,
$\kappa \in \{\mathrm{RGSS}\} \cup \{\mathrm{Span}(a,b)\}$ the span,
$r$ the number of purification rounds applied, and $h \in
\{\mathrm{pending}, \mathrm{resolved}\}$ the heralding status.
\TODO{Motivate each field and state the \emph{legality constraint}: two
states may be purified together only if they share the same $\kappa$.}

\section{Schedule as a Formal Object}
\label{sec:model:schedule}

\subsection{DAG Representation}
\label{sec:model:dag}

\TODO{Define a schedule $\Sigma = (T, \phi)$ as a directed acyclic graph
whose nodes are typed operations
(\texttt{Gen}, \texttt{AbsaBsm}, \texttt{Join}, \texttt{Purify},
\texttt{Herald}, \texttt{PauliCorrect}, and \texttt{Idle}) and whose
edges encode data dependencies. Explain the semantics and parameters of
each node type. There is no \texttt{Discard} node in the implemented
schedule model.}

\subsection{Canonical Timing Scenarios}
\label{sec:model:timing}

\TODO{Reproduce the three closed-form generation- and memory-time
expressions from the Integrating paper (raw, baseline, optimistic) and
show they are special cases of the general DAG evaluator.}

\section{Objective Functions and Resource Budget}
\label{sec:model:objectives}

\TODO{Define the implemented resource budget $e_{\max}$ as the maximum
number of \texttt{GenNode} leaves (half-RGS resources) consumed per trial.
Distinguish it from $M_{\max}$, the maximum concurrent open branches,
which is represented in the resource model but is not enforced by the
search or DAG validator. State the two candidate objectives:
(1)~fidelity subject to a throughput floor, and
(2)~throughput subject to a fidelity floor $F \ge 0.9$.}

\section{Optimality Scope and Known Limitations}
\label{sec:model:scope}

\TODO{State precisely what the recursion finds natively: per-hop
link-level purification with variable copies and circuits, arbitrary
span-partition joins, and one two-copy same-span pumping move. Explain
that pumping candidates and stored frontiers are capped by default, so
default DP and beam results are heuristic. \texttt{exact\_pumping=True}
removes these caps but is only practical for very small instances.
Reference the merged fixed families and the remaining non-guarantee:
the search need not find every optimal legal schedule.}